\begin{titlepage}
\thispagestyle{empty}
\sffamily
\vspace*{1.5cm}
{\color{orange}\large\bfseries RAPPORT SCIENTIFIQUE, PÉDAGOGIQUE ET PROSPECTIF\par}
\vspace{0.8cm}
{\color{bleu}\fontsize{28}{34}\selectfont\bfseries L'économie française\\[2pt] au banc d'essai\par}
\vspace{0.5cm}
{\color{bleu}\Large Ce que montre un modèle ouvert de la France :\\ mécanismes, expériences, stratégies de redressement\par}
\vspace{0.8cm}
{\color{bleu}\rule{\linewidth}{1.2pt}}
\vspace{0.6cm}

{\large\color{gris}Projet « Présidentielle » — modèle économique et social de la France, version 0.6\par}
\vspace{1.2cm}
\renewcommand{\arraystretch}{1.5}
\begin{tabular}{@{}>{\bfseries\color{bleu}}p{4cm}p{10.5cm}@{}}
Date & 18 novembre 2026 \\
Version du modèle & 0.6 (code et résultats : dépôt GitHub du projet, dossier \texttt{rapport/}) \\
Données & Exclusivement ouvertes : Insee, Eurostat, DGFiP, Drees, Banque de France, FMI, COR, douanes, SDES \\
Statut & Document de travail scientifique. Les résultats sont conditionnels aux hypothèses du modèle, qui sont toutes exposées. Ils ne constituent ni une prévision, ni une recommandation politique. \\
\end{tabular}
\vfill
{\small\color{gris}Rédaction et calculs : Claude (Anthropic), à la demande et sous la direction du porteur du projet. Toutes les simulations citées ont été exécutées avec le code du dépôt ; les fichiers de résultats sont fournis avec ce rapport.\par}
\end{titlepage}
